\documentclass[11pt]{article}

\usepackage[margin=1in]{geometry}
\usepackage[numbers]{natbib}


\usepackage{amsmath,amssymb,amsthm}
\usepackage{enumitem}
\usepackage{mathrsfs}
\usepackage{xcolor}
\usepackage{mathtools}
\usepackage{hyperref}
\usepackage[normalem]{ulem}

\usepackage{tikz}
\usepackage{tikz-cd}
\usetikzlibrary{arrows.meta, positioning, cd}
\usepackage{float}

\newtheorem{definition}{Definition}
\newtheorem{assumption}{Assumption}
\newtheorem{lemma}{Lemma}
\newtheorem{proposition}{Proposition}
\newtheorem{theorem}{Theorem}
\newtheorem{remark}{Remark}
\newtheorem{corollary}{Corollary}

\newcommand{\Rcal}{\mathcal{R}}
\newcommand{\Pcal}{\mathcal{P}}
\newcommand{\Mcal}{\mathcal{M}}
\newcommand{\Ncal}{\mathcal{N}}
\newcommand{\Fcal}{\mathcal{F}}
\newcommand{\Bcal}{\mathcal{B}}
\newcommand{\Xcal}{\mathcal{X}}
\newcommand{\Dcal}{\mathcal{D}}

\newcommand{\Psf}{\mathsf{P}}
\newcommand{\Lsf}{\mathsf{L}}
\newcommand{\Rsf}{\mathsf{R}}

\newcommand{\push}{\sharp}     	% pushforward
\newcommand{\Ext}{\mathrm{Ext}} 	% extension

\newcommand{\R}{\mathbb{R}}
\newcommand{\Pbb}{\mathbb{P}}
\newcommand{\Ebb}{\mathbb{E}} 	% expectation 
\newcommand{\Fib}{\mathrm{Fib}} 	% compatibility fibre


\pagecolor{white}


\begin{document}

\tableofcontents

\section{Introduction}

We work throughout on standard Borel spaces, ensuring the existence of regular conditional
probabilities and allowing conditional distributions to be represented as Markov kernels when
such disintegrations exist.

\subsection{Empirical dynamics under boundary queries}

\paragraph{Observable regularities and query families.}
We begin from a class $\mathcal Q$ of \emph{boundary queries}.  
A boundary query is an operational procedure that extracts observable structure from an
underlying stream of observations.  Formally, each query $Q \in \mathcal Q$ is associated with
an outcome space $\mathcal O_Q$ and is understood as a measurable projection, functional, or
disintegration acting on the observable stream.

For each query $Q \in \mathcal Q$, we specify a set of admissible regularities
\[
\mathcal R_Q \subseteq \mathcal P(\mathcal O_Q),
\]
representing the class of empirically stable statistical structures that may arise under that
query.

A \emph{boundary horizon (query-relative)} is a family
\[
\{\nu_Q\}_{Q \in \mathcal Q}, \qquad \nu_Q \in \mathcal R_Q,
\]
interpreted as the empirically observed collection of regularities across all admissible
queries.

\paragraph{Completion as an optional representation.}
When the observable family $\{\nu_Q\}_{Q\in\mathcal Q}$ is jointly realizable, there exists a
path-space measure $\mathbb P_Y \in \mathcal P(Y^{\mathbb Z})$ such that
\[
\nu_Q = Q_\# \mathbb P_Y
\qquad \text{for all } Q \in \mathcal Q.
\]
In this case, $\mathbb P_Y$ provides a \emph{completion} or \emph{process-level representation}
of the observable horizon.  Importantly, the existence of such a completion is not assumed.
Instead, it is treated as an empirical and mathematical question.

Thus, throughout, the observable family $\{\nu_Q\}$ is taken as primitive, while path laws
are regarded as justified when coherence across queries permits.

\paragraph{Justified dynamical structure.}
Finite-resolution and directional dynamical objects (regularities, conditional kernels,
operators) arise from statistical interrogation of the boundary horizon.  Concretely, such
objects are obtained through projection, conditioning, or coarse-graining relative to auxiliary
choices (such as a lag $\Delta t$ or a conditioning $\sigma$-field).  In this sense, dynamical
descriptions are always relative to a querying regime.

When a completion $\mathbb P_Y$ exists, these objects may be realized concretely as measurable
operations on the path space $Y^{\mathbb Z}$ evaluated under $\mathbb P_Y$.  Otherwise, they are
understood operationally through their action on the observable family $\{\nu_Q\}$.

\paragraph{Stationarity and scope.}
Stationarity is invoked only when global index-shift invariance or semigroup structure is
required.  The notions of boundary queries, compatibility, and (approximate) boundary
completion depend only on local extendability of the observable regularities
$\{\nu_Q\}$ across resolutions.  Failure of stationarity does not preclude meaningful
multi-resolution structure, but it may obstruct the identification of a single
time-homogeneous generator.

Failure of (exact or approximate) boundary completion is itself a meaningful empirical outcome.
Such failure indicates that the observable regularities do not admit a coherent multi-resolution
closure and therefore identifies regimes in which new hypotheses or refined queries may be
required.

\paragraph{Boundary queries and query-relative compatibility.}
A boundary query acts directly on the observable horizon.  When a completion exists, it may be
realized as a measurable map or Markov kernel
\[
Q : Y^{\mathbb Z} \leadsto \mathcal O_Q,
\]
inducing the law $Q_\# \mathbb P_Y$.  More generally, compatibility is defined operationally.

Let $\mathcal Q$ be a fixed class of admissible queries.  A representation
$s$ is said to be \emph{$\mathcal Q$-compatible} if each observable query can be reproduced after
passing through the representation:
\[
s \text{ is } \mathcal Q\text{-compatible}
\quad \Longleftrightarrow \quad
\forall Q \in \mathcal Q \; \exists \widetilde Q \text{ such that } Q = \widetilde Q \circ s.
\]
This expresses the idea that the representation preserves all observable structure specified by
the query class.

\paragraph{Finite-lag and structural queries.}
No single query is privileged.  Instead, we consider families of finite-resolution and
structural queries adapted to the empirical context.

A basic example consists of finite-lag projections.  For each $\Delta t \in H$, define
\[
Q_{\Delta t} : Y^{\mathbb Z} \to Y \times Y,
\qquad
Q_{\Delta t}(y) = (y_0, y_{\Delta t}),
\]
which induce the corresponding lag-regularities
\[
\Gamma_Y^{(\Delta t)} := (Q_{\Delta t})_\# \mathbb P_Y
\]
when a completion exists.  These summaries serve as the primary objects for
cross-resolution comparison and extension.

More general query families may include multi-time projections, coarse-grainings, or
experimentally defined observables.  Compatibility and completion are therefore always
understood relative to the chosen class $\mathcal Q$.

\paragraph{Conditional and predictive structure as derived queries.}
Conditional or predictive descriptions arise as special cases when suitable disintegrations
exist.  For example, conditioning future observations on the past yields predictive kernels.
Within the present framework, such structures are treated as derived statistical queries rather
than as primitive dynamical assumptions.  Their relevance depends on the empirical setting and
the class of admissible queries.

% ============================================================

\section{Finite-lag objects as boundary queries}

We now introduce the basic observable summaries obtained by applying
finite-resolution queries to the empirical horizon. These summaries form the
primary objects used for comparison across observational scales.

\subsection{Edge regularities as primary finite-lag objects}

Let $Y$ be a
standard Borel space $Y$, and let $\pi_t:Y^{\mathbb Z}\to Y$ denote the
coordinate projection at time $t$.

\begin{definition}[Edge regularity]\label{def:edge-regularity}
For any lag $\Delta t\in H$, the \emph{edge regularity} at scale $\Delta t$ is the two-time
regularity $\Gamma_Y^{(\Delta t)}\in\mathcal P(Y\times Y)$ induced by the lag query
$Q_{\Delta t}(y)=(y_0,y_{\Delta t})$.
When a path-law completion $\mathbb P_Y\in\mathcal P(Y^{\mathbb Z})$ exists, this is the
pushforward $(\pi_0,\pi_{\Delta t})_\#\mathbb P_Y$.
\end{definition}

We interpret $\Gamma_Y^{(\Delta t)}$ as the outcome of a \emph{finite-lag
query}---an observational procedure, not a dynamical law.  Different choices of
$\Delta t$ correspond to different resolutions at which the empirical horizon is
interrogated.

When marginal coherence holds across the family $\{\Gamma_Y^{(\Delta t)}\}$ (see
Section~\ref{sec:boundary-completion}), we denote the common one-time marginal by
\[
\mu_Y\in\mathcal P(Y),
\]
and note that $\mu_Y=(\pi_0)_\#\mathbb P_Y$ when a completion exists.

\subsection{Multi-time regularities}

More detailed summaries arise from projections onto multiple coordinates.

\begin{definition}[Three-time regularity]\label{def:three-time}
For lags $\Delta t_1,\Delta t_2\in H$, define the three-time regularity
\[
\Gamma_Y^{(\Delta t_1,\Delta t_2)}
:=
(\pi_0,\pi_{\Delta t_1},\pi_{\Delta t_1+\Delta t_2})_\#\,\mathbb P_Y
\;\in\; \mathcal P(Y^3).
\]
\end{definition}

These summaries capture observable relationships across multiple scales and
will play a central role in assessing cross-resolution coherence.

\subsection{Kernel representations (optional)}

Whenever an edge regularity admits a disintegration, one may introduce a
kernel representation.

\begin{definition}[Kernel representation]\label{def:kernel-edge}
If
\[
\Gamma_Y^{(\Delta t)}(dy,dy')
=
\mu_Y(dy)\,T_{\Delta t}(y,dy'),
\]
then $T_{\Delta t}:Y\leadsto Y$ is called a kernel representation of the
lag-$\Delta t$ regularity.
\end{definition}

Such representations provide convenient operator forms but are not required.
Given a kernel, one obtains a linear operator on observables $f:Y\to\mathbb R$:
\[
(T_{\Delta t}f)(y)
:=
\int_Y f(y')\,T_{\Delta t}(y,dy').
\]

\subsection{Markov completion (derived)}

A family of kernels $\{T_{\Delta t}\}_{\Delta t\in H}$ is said to form a
Markov completion when it approximately closes across lags:
\[
T_{\Delta t_1+\Delta t_2}
\;\approx\;
T_{\Delta t_1}\circ T_{\Delta t_2},
\]
with approximation quantified by a divergence criterion on the induced
edge regularities.

Continuity in $\Delta t$ may then permit the introduction of a generator as
a limiting representation.

% ============================================================

\section{Boundary horizons and query-relative compatibility}

\subsection{Boundary primacy and non-privileging}

\begin{proposition}[Boundary primacy]\label{prop:boundary-primacy}
Let $Y$ be a standard Borel space and let $\mathbb P_Y \in \mathcal P(Y^{\mathbb Z})$
denote a boundary horizon. All finite-resolution and directional structures
(e.g.\ regularities, kernels, operators) arise relative to a boundary query.
No finite resolution is intrinsically distinguished.
\end{proposition}

\begin{proof}
This follows from the operational construction of observable structure
introduced in Section~1.
\end{proof}

\begin{remark}[Non-privileging of resolution]
No canonical temporal discretization, lag, or compositional scale is assumed.
Observable summaries depend on the chosen querying regime.
\end{remark}

\subsection{Boundary queries}

\begin{definition}[Boundary query]\label{def:boundary-query}
A boundary query is any measurable operation on the path space
realized by projection, disintegration, or composition thereof.
Equivalently, it may be represented as a Markov kernel
\[
Q:Y^{\mathbb Z}\leadsto \mathcal O,
\]
inducing a regularity $Q_\#\mathbb P_Y$.
\end{definition}

\subsection{Query-relative compatibility}

Let $\mathcal Q$ be a class of admissible queries.

\begin{definition}[Compatibility]\label{def:query-compatibility}
A representation
\[
s:Y^{\mathbb Z}\to \mathcal S
\]
is $\mathcal Q$-compatible if for every $Q\in\mathcal Q$ there exists
$\widetilde Q$ such that
\[
Q=\widetilde Q\circ s.
\]
\end{definition}

\begin{proposition}[Universal property]\label{prop:compat-universal}
A representation is compatible if and only if all queries in $\mathcal Q$
factor through it.
\end{proposition}

\begin{proof}
Immediate from Definition~\ref{def:query-compatibility}.
\end{proof}

\begin{remark}
Compatibility is relative to $\mathcal Q$; distinct compatible representations
may differ substantially in internal structure.
\end{remark}

% ============================================================

\section{Boundary completion across resolutions}
\label{sec:boundary-completion}

We now study when observable regularities at different resolutions exhibit coherent
cross-scale structure.  The notion of \emph{boundary completion} is empirical and
pre-dynamical: it concerns joint extendability (consistency) of query-induced
regularities across resolutions, and does not assume any underlying evolution law.

\subsection{Boundary queries and derived regularities}
\label{subsec:boundary-queries-derived}

Fix a set of lags $H\subset\mathbb R_+$ (or $H\subset\mathbb Z_+$).  For each $\Delta t\in H$,
consider the finite-lag query
\[
Q_{\Delta t}:Y^{\mathbb Z}\to Y\times Y,
\qquad
Q_{\Delta t}(y)=(y_0,y_{\Delta t}).
\]
A query-relative horizon specifies the induced two-time regularity at lag $\Delta t$ as
\[
\Gamma_Y^{(\Delta t)} \in \mathcal P(Y\times Y).
\]
When a path-law completion $\mathbb P_Y\in\mathcal P(Y^{\mathbb Z})$ exists and $Q_{\Delta t}$ is
realized as a measurable projection, this agrees with the pushforward
$\Gamma_Y^{(\Delta t)}=(Q_{\Delta t})_\#\mathbb P_Y$.

We write
\[
\mathcal B_Y(H):=\{\Gamma_Y^{(\Delta t)}:\Delta t\in H\}\subset\mathcal P(Y\times Y)
\]
for the family of lag-$H$ boundary regularities.

\subsection{Projective compatibility and boundary completion}
\label{subsec:projective-compatibility}

We now make precise the relationship between query-relative boundary regularities and
classical consistency conditions for stochastic processes.  The key point is that the
boundary framework generalizes Kolmogorov consistency from coordinate projections to
arbitrary classes of observational queries.

\paragraph{Query systems as projective systems.}
Let $\mathcal Q$ be a class of boundary queries, each associated with a standard Borel
outcome space $\mathcal O_Q$.  Assume $\mathcal Q$ is equipped with a preorder $\preceq$
expressing refinement: for $Q_1,Q_2\in\mathcal Q$, we write $Q_1\preceq Q_2$ if $Q_2$ refines
$Q_1$, meaning that there exists a measurable map
\[
\pi_{21}:\mathcal O_{Q_2}\to\mathcal O_{Q_1}
\]
such that
\[
Q_1=\pi_{21}\circ Q_2.
\]
Thus the query class induces a projective system of measurable spaces.

\paragraph{Projective families of boundary regularities.}
A family of boundary regularities
\[
\{\nu_Q\}_{Q\in\mathcal Q},
\qquad
\nu_Q\in\mathcal P(\mathcal O_Q),
\]
is said to be \emph{projectively compatible} if for all $Q_1\preceq Q_2$,
\[
\nu_{Q_1}=(\pi_{21})_\#\nu_{Q_2}.
\]
This condition expresses exact stability of observable regularities under refinement of the
querying regime.

\paragraph{Generalized Kolmogorov extension.}
\begin{theorem}[Generalized Kolmogorov extension]\label{thm:generalized-kolmogorov}
Let $Y$ be a standard Borel space and let $\mathcal Q$ be a directed class of boundary
queries with refinement preorder.  Suppose that the family $\{\nu_Q\}_{Q\in\mathcal Q}$ is
projectively compatible and that the $\sigma$-algebra generated by the query cylinders
separates points in $Y^{\mathbb Z}$.  Then there exists a probability measure
\[
\mathbb P_Y\in\mathcal P(Y^{\mathbb Z})
\]
such that
\[
\nu_Q = Q_\#\mathbb P_Y
\qquad \text{for all } Q\in\mathcal Q.
\]
Moreover, this completion is unique on the $\sigma$-algebra generated by the query family.
\end{theorem}

\begin{remark}[Classical Kolmogorov as a special case]
If $\mathcal Q$ consists of all finite coordinate projections, then projective compatibility
reduces to classical Kolmogorov consistency.  Thus stochastic processes arise as special cases
of compatibility across coordinate queries.
\end{remark}

\begin{remark}[Interpretation]
The theorem shows that a path-space boundary horizon, when it exists, can be regarded as the
projective limit of observable regularities.  In this sense, global path-law structure is not
assumed but arises as a completion of compatible empirical summaries.
\end{remark}

\subsection{Exact boundary completion}
\label{subsec:exact-boundary-completion}

\begin{definition}[Exact completion]\label{def:exact-boundary-completion}
The family $\mathcal B_Y(H)$ is \emph{boundary complete} if:
\begin{enumerate}
  \item (\emph{Marginal coherence}) All lag regularities share a common one-time marginal:
  \[
  (\pi_0)_\#\Gamma_Y^{(\Delta t)}=\mu_Y
  \qquad\text{for all }\Delta t\in H,
  \]
  for some $\mu_Y\in\mathcal P(Y)$.
  \item (\emph{Two-step extendability}) For all $\Delta t_1,\Delta t_2\in H$ with
  $\Delta t_1+\Delta t_2\in H$, the extension set
  \[
  \Ext(\Gamma_Y^{(\Delta t_1)},\Gamma_Y^{(\Delta t_2)})
  \]
  is nonempty.
\end{enumerate}
\end{definition}

This condition expresses joint extendability of observable summaries, and precedes any Markov
or semigroup interpretation.

\subsection{Approximate completion}
\label{subsec:approx-boundary-completion}

Exact consistency is often too strong. We therefore quantify deviations from extendability.

\begin{definition}[$\varepsilon$--completion]\label{def:eps-boundary-completion}
The family is $\varepsilon$--boundary complete if for all $\Delta t_1,\Delta t_2\in H$ with
$\Delta t_1+\Delta t_2\in H$,
\[
\inf_{\widetilde\Gamma\in\Ext(\Gamma_Y^{(\Delta t_1)},\Gamma_Y^{(\Delta t_2)})}
D_f\!\bigl(\Gamma_Y^{(\Delta t_1+\Delta t_2)},\widetilde\Gamma\bigr)
\le \varepsilon.
\]
\end{definition}

\subsection{Derived dynamics and earned structure}

Boundary completion determines which additional structures may be introduced
without contradicting observable regularities.

\begin{itemize}
  \item Kernel representations arise only when edge regularities admit
  disintegration.

  \item Semigroup structure reflects cross-resolution closure of the family
  $\mathcal B_Y(H)$.

  \item Generators, when they exist, arise as limiting objects associated with
  increasingly refined observational scales.
\end{itemize}

Interior dynamical realizations are therefore admissible only insofar as they
reproduce the completed boundary structure.  They serve as explanatory
devices rather than primitive sources of empirical regularity.

\medskip

In this sense, all dynamical structure in the framework is \emph{earned from the boundary}, rather
than imposed upon it.
% ============================================================


\section{Lag queries, sufficiency, and dominance}

\subsection{Finite-lag boundary queries}

Throughout, we treat the family of two-time boundary regularities
\[
\mathcal B_Y(H)
:= \bigl\{\Gamma_Y^{(\Delta t)} \in \mathcal P(Y\times Y)\bigr\}_{\Delta t\in H}
\]
as the primary finite-resolution empirical summaries of the boundary horizon.
When a path-space regularity $\mathbb P_Y\in\mathcal P(Y^{\mathbb Z})$ exists,
these may be realized as
\[
\Gamma_Y^{(\Delta t)} = (\pi_0,\pi_{\Delta t})_\#\mathbb P_Y,
\]
but we do \emph{not} assume such a completion a priori.

We refer to each $\Gamma_Y^{(\Delta t)}$ as a \emph{lag-$\Delta t$ boundary query}.
Such queries are understood as statistical observations of the boundary horizon,
not as mechanisms that generate it.

\begin{remark}[No privileged time step]
We do not single out any distinguished ``one-step'' lag or canonical discretization.
All $\Delta t\in H$ are treated symmetrically; structure is only admissible under relations
among these queries.
\end{remark}

\subsection{Lag sufficiency}

\begin{definition}[Lag sufficiency]\label{def:lag-sufficiency}
The boundary family $\mathcal B_Y(H)$ is said to be \emph{lag-sufficient} if every
finite-dimensional boundary regularity that is empirically meaningful for the system
is determined (up to the ambiguity of extension) by the collection
$\{\Gamma_Y^{(\Delta t)}\}_{\Delta t\in H}$.
\end{definition}

Intuitively, lag sufficiency expresses that all higher-order finite-resolution
summaries are shadows of, or constraints upon, the two-time regularities.

\begin{proposition}[Lag dominance]\label{prop:lag-dominance}
Suppose $\mathcal B_Y(H)$ is exactly boundary complete (Definition~\ref{def:exact-boundary-completion}).
Then every finite-dimensional marginal of any admissible path-space completion of the
boundary horizon is compatible with, and constrained by, the family
$\{\Gamma_Y^{(\Delta t)}\}_{\Delta t\in H}$.
\end{proposition}

\begin{proof}[Sketch]
Exact boundary completion ensures that for any finite collection of times
$t_0<t_1<\cdots<t_n$ with increments in $H$, there exists a joint law on
$Y^{n+1}$ whose adjacent marginals coincide with the corresponding edge regularities.
Thus any higher-order finite-dimensional summary must be realizable via successive
extensions of lag queries, and hence cannot contradict them.
\end{proof}

\begin{corollary}[Minimality of lag sufficiency]\label{cor:lag-minimality}
Any representation that preserves all lag-$\Delta t$ boundary regularities
$\{\Gamma_Y^{(\Delta t)}\}_{\Delta t\in H}$ preserves all finite-lag empirical content
of the boundary horizon.
\end{corollary}

\begin{remark}[Lag sufficiency is not Markovity]\label{rem:lag-not-markov}
Lag sufficiency does \emph{not} imply the existence of Markov kernels, conditional
independence, or a semigroup structure. It is a statement about empirical summaries,
not about dynamics.
\end{remark}

\subsection{Relation to predictive conditioning (optional)}

If a path-space regularity $\mathbb P_Y$ exists and admits disintegration, one may
form the predictive kernel
\[
K_\infty(\cdot\mid p)
:= \mathbb P_Y(Y_{0:\infty}\in\cdot \mid Y_{-\infty:0}=p).
\]
In that case, each $\Gamma_Y^{(\Delta t)}$ is a marginal of $K_\infty$.

However, we do \emph{not} take $K_\infty$ to be primitive or canonical. Rather, it is
a derived object that may be useful in some contexts, but whose relevance is contingent
on additional structure beyond lag sufficiency and boundary completion.



% ============================================================

\section{Extension and boundary completion}

\subsection{Extension as joint realizability}

\begin{definition}[Extension set]\label{def:ext-set}
Let $\Gamma_1,\Gamma_2\in\mathcal P(Y\times Y)$ be two boundary regularities for successive lags
(e.g.\ $\Gamma_1=\Gamma_Y^{(\Delta t_1)}$ and $\Gamma_2=\Gamma_Y^{(\Delta t_2)}$).
Assume they share a compatible intermediate marginal in the sense that there exists
$\mu\in\mathcal P(Y)$ such that
\[
(\pi_2)_\#\Gamma_1=\mu=(\pi_1)_\#\Gamma_2.
\]
Define the \emph{extension set}
\[
\Ext(\Gamma_1,\Gamma_2)
:=
\Bigl\{(\pi_0,\pi_2)_\#\Lambda :
\Lambda\in\mathcal P(Y^3),\
(\pi_0,\pi_1)_\#\Lambda=\Gamma_1,\
(\pi_1,\pi_2)_\#\Lambda=\Gamma_2
\Bigr\}.
\]
Elements of $\Ext(\Gamma_1,\Gamma_2)$ are precisely the endpoint regularities that can arise from
a three-time joint law consistent with the prescribed two-time regularities.
\end{definition}

\begin{proposition}[Existence of extensions]\label{prop:ext-existence}
$\Ext(\Gamma_1,\Gamma_2)\neq\varnothing$ if and only if $\Gamma_1$ and $\Gamma_2$ share a common
intermediate marginal $\mu$.
\end{proposition}

\begin{proof}[Sketch]
If $\Ext(\Gamma_1,\Gamma_2)\neq\varnothing$, any witness $\Lambda$ must satisfy
$(\pi_2)_\#\Gamma_1=(\pi_1)_\#\Gamma_2$ by construction.
Conversely, if such a marginal $\mu$ exists, standard gluing arguments for probability measures
on standard Borel spaces guarantee the existence of some $\Lambda$ with the required pairwise
marginals.
\end{proof}

\begin{remark}[Extension is not composition]\label{rem:ext-not-comp}
$\Ext$ expresses a condition of \emph{joint realizability} among empirical regularities.
It does not assume Markov kernels, conditional independence, causal direction, or any
compositional law on $Y$. Time order and dynamics are not built into the definition.
\end{remark}

\subsection{Approximate boundary completion}

Fix a divergence $D_f(\cdot,\cdot)$ on probability measures and a tolerance $\varepsilon>0$.

\begin{definition}[$\varepsilon$--boundary completion]\label{def:eps-boundary-completion}
The boundary family $\mathcal B_Y(H)$ is said to be \emph{$\varepsilon$--boundary complete} if for all
$\Delta t_1,\Delta t_2\in H$,
\[
\inf_{\widetilde\Gamma\in
\Ext\bigl(\Gamma_Y^{(\Delta t_1)},\Gamma_Y^{(\Delta t_2)}\bigr)}
D_f\!\Big(
(\pi_0,\pi_{\Delta t_1+\Delta t_2})_\# \mathbb P_Y,\;
\widetilde\Gamma
\Big)
\;\le\;
\varepsilon.
\]
\end{definition}

Conceptually, $\varepsilon$--boundary completion measures how closely the empirically observed
long-lag boundary regularity can be approximated by \emph{some} joint realization consistent with
the corresponding shorter-lag regularities.

This definition refers only to joint extendability of boundary summaries and does \emph{not}
assume Markov kernels, conditional independence, or causal composition.

\begin{remark}[Interpretation of $\varepsilon$]\label{rem:epsilon-meaning}
The parameter $\varepsilon$ admits three complementary readings:
\begin{enumerate}
  \item As a measure of \emph{non-extendability} across resolutions;
  \item As a bound on \emph{mutual expressibility} of boundary viewpoints;
  \item As an \emph{operational budget} for hypothesis generation when exact completion fails.
\end{enumerate}
Large $\varepsilon$ signals genuine incompatibility of boundary summaries, not merely data noise.
\end{remark}

% ============================================================

\section{Admissible dynamical structure}

The preceding sections define boundary regularities and their cross-resolution consistency
without invoking any notion of dynamics. We now describe how familiar dynamical objects—Markov
kernels, semigroups, and generators—may be \emph{admissible} as representations of boundary structure
when additional regularity is present. Nothing in this section is assumed at the boundary level;
all constructions are conditional and generally non-unique.

\subsection{From extension to kernel representations}

\begin{proposition}[Admissible kernel representation]\label{prop:admissible-kernel}
Let $\mathcal B_Y(H)$ be an $\varepsilon$--boundary complete family of boundary regularities.
Suppose that for each $\Delta t \in H$, the corresponding edge regularity
$\Gamma_Y^{(\Delta t)} \in \mathcal P(Y\times Y)$ admits a disintegration of the form
\[
\Gamma_Y^{(\Delta t)}(dy,dy')
=
\mu_Y(dy)\,T_{\Delta t}(y,dy'),
\]
for some probability kernel $T_{\Delta t}:Y\leadsto Y$.
Then $\{T_{\Delta t}\}_{\Delta t\in H}$ provides a kernel-based representation of the boundary
regularities consistent with the given querying regime.
\end{proposition}

\begin{remark}
The existence of kernels $T_{\Delta t}$ is not guaranteed by boundary completion alone.
Disintegration is an additional structural property of the boundary regularities, and kernel
representations are introduced only when such disintegrations exist.
\end{remark}

\begin{remark}[Non-uniqueness of kernel representations]\label{rem:kernel-nonunique}
Even when disintegrations exist, kernel representations are generally non-unique.
Multiple families of kernels may induce the same boundary regularities and hence be equally
compatible with the boundary horizon. This reflects underdetermination at the level of interior
dynamics.
\end{remark}

\subsection{Approximate semigroup structure}

\begin{proposition}[Approximate semigroup emergence]\label{prop:approx-semigroup}
Suppose that $\mathcal B_Y(H)$ is $\varepsilon$--boundary complete and admits kernel
representations $\{T_{\Delta t}\}_{\Delta t\in H}$. Then for $\Delta t_1,\Delta t_2\in H$,
the composed kernel $T_{\Delta t_1}\circ T_{\Delta t_2}$ approximates $T_{\Delta t_1+\Delta t_2}$
in the sense that
\[
D_f\!\Big(
\mu_Y \otimes (T_{\Delta t_1}\circ T_{\Delta t_2}),\;
\Gamma_Y^{(\Delta t_1+\Delta t_2)}
\Big)
\;\lesssim\;
\varepsilon,
\]
where the comparison is mediated by the boundary completion tolerance.
\end{proposition}

\begin{proof}[Sketch]
By $\varepsilon$--boundary completion, there exists some
$\widetilde\Gamma\in
\Ext(\Gamma_Y^{(\Delta t_1)},\Gamma_Y^{(\Delta t_2)})$
such that
\[
D_f\!\Big(
\Gamma_Y^{(\Delta t_1+\Delta t_2)},\widetilde\Gamma
\Big)\le \varepsilon.
\]
When disintegrations exist, the product-form endpoint law
$\mu_Y\otimes(T_{\Delta t_1}\circ T_{\Delta t_2})$ is an admissible element of the same extension
set. Selecting this witness yields the stated approximation.
\end{proof}

\begin{remark}
The approximate semigroup property is not imposed but follows from the consistency of boundary
regularities across resolutions. The quality of the approximation is controlled by $\varepsilon$
and may vary across scales.
\end{remark}

\subsection{Generators as closure limits}

\begin{theorem}[Generator as closure limit]\label{thm:generator-limit}
Assume that the family $\{T_{\Delta t}\}_{\Delta t\in H}$ admits a limit in an appropriate operator
topology as $\Delta t\to 0$. When such a limit exists, it defines an infinitesimal generator $L$
that provides a continuous-time dynamical representation of the boundary regularities.
\end{theorem}

\begin{remark}[Conditional status of generators]
The existence, uniqueness, and interpretation of generators depend on additional analytical
assumptions beyond boundary completion, including continuity, topology, and limiting behavior.
Generators are therefore optional representational tools rather than mandatory consequences of
the framework.
\end{remark}

\subsection{Empirical adequacy versus explanatory realization}

\begin{remark}[Boundary adequacy and explanation]\label{rem:adequacy-vs-explain}
Boundary compatibility and completion constrain empirical adequacy.
Kernel-based dynamics and generators provide explanatory structure but are not uniquely
determined by boundary data. The framework therefore separates empirical regularity
(boundary adequacy) from explanatory commitment (interior realization).
\end{remark}

\begin{remark}[Over- and under-explanation]
A realization may be empirically adequate yet introduce excess structure not warranted by the
boundary horizon, or conversely may fail to capture boundary regularities despite rich internal
dynamics. Such distinctions become visible only after separating boundary completion from
dynamical representation.
\end{remark}

% ============================================================


\section{Information geometry on boundary queries}

\begin{proposition}[Query-induced geometry]\label{prop:query-geometry}
% Geometry lives on statistical manifolds induced by queries.
\end{proposition}

\begin{theorem}[Invariance under sufficient statistics]\label{thm:cencov-invariance}
% Cite \v{C}encov/Bauer/Ay-type results.
\end{theorem}

\begin{remark}[Divergence choice and commensurability]\label{rem:divergence-choice}
% Divergences as empirical interfaces, not mere tolerances.
\end{remark}


% ============================================================


\section{Underdetermination and empirical stance}

\begin{proposition}[Multiplicity of compatible realizations]\label{prop:many-realizations}
% Boundary adequacy does not fix interior explanation.
\end{proposition}

\begin{remark}[Empirical adequacy vs explanation]\label{rem:adequacy-vs-explain}
% Clear separation of roles.
\end{remark}

\begin{remark}[Failure of completion as signal]\label{rem:completion-failure}
% Legitimate empirical outcome.
\end{remark}

\subsection*{Reflection questions}

The following questions are intended to audit the conceptual stance of the framework and to
identify remaining points of pressure or underdetermination. They are not requirements imposed
by the formalism, but diagnostics for its interpretation and scope.

\paragraph{Boundary completion and empirical meaning.}
\begin{enumerate}
  \item Does the definition of exact boundary completion successfully characterize
  \emph{empirical coherence} among finite-resolution regularities without presupposing any
  underlying generative or causal process?

  \item Can failure of boundary completion be an informative empirical outcome---i.e.\ does it
  distinguish genuine incompatibility of regularities from insufficient data or estimation error?

  \item Is approximate boundary completion ($\varepsilon>0$) clearly interpreted as a graded
  empirical property of observed regularities, rather than as a defect to be repaired by modeling?
\end{enumerate}

\paragraph{Extension versus dynamics.}
\begin{enumerate}
  \item Is the notion of extension (\(\Ext\)) unambiguously understood as a condition of
  \emph{joint realizability}, rather than as an implicit form of temporal composition or causal
  evolution?

  \item After Sections~6 and~7, is it clear that kernel composition and semigroup structure are
  representational conveniences that may be admissible, rather than assumptions imposed at the
  boundary level?

  \item Would a reader unfamiliar with Markov processes nonetheless understand extension as a
  purely probabilistic constraint, independent of time direction or mechanism?
\end{enumerate}

\paragraph{Admissible structure and underdetermination.}
\begin{enumerate}
  \item Is it explicit---and acceptable within the framework---that multiple, non-equivalent kernel
  families or generators may be equally compatible with the same boundary regularities?

  \item Is the separation between \emph{boundary adequacy} (empirical compatibility) and
  \emph{explanatory realization} (interior dynamical structure) sufficiently sharp to prevent
  conflating predictive success with causal explanation?

  \item When multiple realizations are boundary-compatible, does the framework leave their
  selection as an external (pragmatic or scientific) choice rather than an internal requirement?
\end{enumerate}

\paragraph{Role and interpretation of \(\varepsilon\).}
\begin{enumerate}
  \item Should \(\varepsilon\) be interpreted primarily as a tolerance for empirical noise, a bound
  on cross-resolution expressibility, or an operational budget for hypothesis generation? Is room
  left for all three interpretations?

  \item Are there plausible scenarios in which \(\varepsilon\) is small locally (at some resolutions)
  but large globally, and does the framework accommodate such scale-dependent behavior?
\end{enumerate}

\paragraph{Scope and limits of the stance.}
\begin{enumerate}
  \item Does the framework remain coherent if no boundary horizon exists even as an idealization,
  yet approximate completion holds at all empirically accessible resolutions?

  \item Is it clear which claims in the manuscript are empirical (about observed regularities),
  structural (about joint realizability), and representational (about kernels, semigroups, or
  generators)?

  \item Finally, is this a stance one is willing to be dogmatic about---namely, that all dynamical
  structure must be admissible from boundary regularities and never assumed in advance?
\end{enumerate}
\end{document}